\section{Introduction}\label{sec:introduction}

An Ising perceptron assigns a weight to each configuration
$x\in\{-1,1\}^N$ through its projections onto independent random patterns.
For Gaussian patterns $g^a\in\mathbb R^N$, a single projection is
$g^a\cdot x/\sqrt N$. Given a real log-potential $f$, the configuration
weight is the product of the factors $\exp f(g^a\cdot x/\sqrt N)$.
The logarithm of their uniform average, divided by $N$, measures the exponential scale
of the weighted configuration space. We use the uniform probability prior
on the cube and call this normalized log partition function the pressure.
The number of patterns is proportional to $N$, with fixed density
$\alpha>0$. Section~\ref{sec:model} gives the exact normalization.

This model combines a discrete configuration space with nonlinear random
constraints. At small density or weak interaction, a single typical overlap
can describe the pressure. More generally, a formula must retain the
distribution of overlaps between two samples from the Gibbs measure. The
patterns are Gaussian, but a nonlinear sum of their log-potentials is
generally not a Gaussian process indexed by the spins. Thus the Parisi
formula for Gaussian mixed $p$-spin Hamiltonians does not directly evaluate
this pressure.

\subsection{Storage models and rigorous predecessors}

The statistical-physics study of random-pattern storage was developed by
Gardner~\cite{gardner1988} and Gardner and Derrida~\cite{gardnerderrida1988}.
The latter work explicitly considered Ising couplings alongside spherical
ones. Krauth and M\'ezard~\cite{krauthmezard1989} analyzed the binary-coupling
model and predicted its storage capacity. In a storage problem a pattern
factor may be an indicator that a constraint is satisfied. At positive
temperature, a finite log-potential instead assigns a positive weight to
every configuration, allowing the pressure to measure both the pattern
reward and the number of configurations achieving it.

Talagrand established replica-symmetric formulas in small-density
perceptron regimes~\cite{talagrand2000halfspaces}. For Gaussian Ising
perceptrons with bounded Borel log-potentials, his high-temperature work
proved overlap-fluctuation bounds at sufficiently small
$\alpha$~\cite{talagrand2002}. Bolthausen, Nakajima, Sun, and
Xu~\cite{bolthausen2022gardner} later gave a new small-density proof of the
Gardner formula using first and second moments conditional on approximate
message passing. In their convention the activation is the multiplicative
weight $U=e^f$, up to constant rescaling. Their Theorem~1.1 treats measurable
weights satisfying a Gaussian-reweighted variance bound and
$\mathbb E[G U(G)]\ne0$; Remark~1.2 discusses the complementary zero-moment
case, where a direct second-moment argument applies. Their small-density
conclusions cover weights bounded above and away from zero, in particular
those arising from bounded real $f$. The density restrictions in these predecessors are
substantive and are absent from Theorem~\ref{thm:main}.

Hard constraints have a distinct, extensive theory. Ding and
Sun~\cite{dingsun2025} and Huang~\cite{huang2024capacity} proved matching
lower and upper capacity bounds at zero margin under their respective
explicit analytical conditions. Shmalo~\cite{shmalo2026capacity} subsequently
reported a computer-assisted verification of both conditions, yielding the
predicted zero-margin threshold; that preprint explicitly describes the
verification as a new proof claim not yet independently reproduced or peer
reviewed. Nakajima and Sun~\cite{nakajimasun2022}
proved concentration, sharp threshold sequences, and disorder universality
for broad hard-constraint classes. For symmetric interval constraints,
Perkins and Xu~\cite{perkinsxu2024} obtained entropy concentration and
freezing under an analytical condition; Abbe, Li, and
Sly~\cite{abbeli2021} proved a lognormal partition-function limit and removed
that condition, working with random-sign patterns. These results concern
specific zero-one weights. Such weights correspond to a log-potential
that takes the value $-\infty$, whereas the theorems here concern finite
log-potentials. A hard-constraint limit would require additional control
of the order of limits.

\subsection{The formula and its proof mechanisms}

Theorem~\ref{thm:main} evaluates the limiting Gaussian Ising pressure for
every fixed positive density and every bounded Borel real log-potential.
Writing $f=\beta\phi$ includes every fixed positive inverse temperature
$\beta$ for a bounded activation $\phi$. The variational parameter is a
nondecreasing overlap quantile $q:(0,1)\to[0,1]$. Its pattern contribution
is a Gaussian recursion, and its entropy contribution is the dual of the
one-spin $\log\cosh$ recursion. Section~\ref{sec:model} defines these
functionals before stating the theorem. No symmetry, concavity, or
smallness condition is imposed on the activation.

The proof uses the Ghirlanda--Guerra overlap
identities~\cite{gg1998}, Panchenko's ultrametricity
mechanism~\cite{panchenko2013}, and Ruelle probability
cascades~\cite{ruelle1987}. For Gaussian mixed-spin Hamiltonians,
Guerra's interpolation bound~\cite{guerra2003} led to Talagrand's proof
of the Parisi formula for even interactions~\cite{talagrand2006parisi};
Panchenko~\cite{panchenko2014parisi} subsequently included odd interactions
using ultrametricity.
We supply the array, duplication, and finite-cascade arguments needed
here. In particular, the changes of measure for the Gaussian marks are
proved with the conditioning required later. The joint transformation of
marks and cascade weights was developed by Panchenko and
Talagrand~\cite{panchenkotalagrand2007property,panchenkotalagrand2007interpolation}.

The hierarchical Gaussian-field enrichment is the construction used by
Mourrat and Panchenko~\cite{mourratpanchenko2020} to extend the Parisi
formula along a Hamilton--Jacobi equation. Our upper bound implements
Mourrat's contact-point supersolution strategy for nonconvex
interactions~\cite[Section~4]{mourrat2021}; see also its vector-spin
extension~\cite{mourrat2023upper}. In that strategy, a quadratic penalty
in perturbation parameters gives concentration and overlap identities
at deterministic contact points. Here the interaction parameter is the
Poissonized pattern density, whose derivative adds one nonlinear pattern.
Joint spin and cascade overlap identities, together with the
monotone-coupling argument of
synchronization~\cite{panchenko2015multispecies}, turn the field variations
at a contact point into the tail inequalities needed for this derivative.

The lower bound follows the cavity variational tradition of Aizenman,
Sims, and Starr~\cite{aizenmansimsstarr2003}. Its identification of the
cavity overlap as a derivative of the one-spin recursion is related to
the critical-point description developed by Chen and
Mourrat~\cite{chenmourrat2025}, and to the critical-point bounds of Chen,
Issa, and Mourrat~\cite[Section~6]{chenissamourrat2026}. Those results
concern Gaussian spin-glass Hamiltonians and do not directly evaluate the
nonlinear pattern model. In the present setting the Gaussian cavity
covariance must itself be identified from the bulk model. Two restored
patterns provide both first and second moments for that identification.

The companion article \emph{The free energy of the spherical random
perceptron}~\cite{spherical} develops the corresponding scalar enrichment
and two-pattern cavity strategy for spherical spins. Here all necessary
arguments are included, and the Ising prior requires its own entropy
calculation and cavity identity. The required concavity of the field
functional is the centered-Ising path-convexity theorem proved by Chen,
Issa, and Mourrat~\cite[Proposition~2.2]{chenissamourrat2026} and, by a
Parisi PDE argument, by Ho~\cite[Theorem~1.8]{ho2026concavity}, with the
opposite sign and a change of covariance normalization. We include a proof
covering repeated field levels and zero initial variance. Coordinate
exchangeability then identifies the cavity-spin overlap with the bulk
overlap. In the cascade calculation this identifies the dual subgradient;
the cited concavity turns this identity into attainment of the entropy
supremum. The model-specific task is to establish these identities for
the nonlinear perceptron.

\subsection{Organization of the proof}

The proof first treats smooth compactly supported log-potentials.
The two bounds use the variational paths differently. The upper comparison
starts from an arbitrary trial overlap path. The lower bound must instead
use the overlap path and field generated by the bulk model, and prove that
this field attains the entropy dual at that same overlap path.
Section~\ref{var:section} proves first variations, the tail order of the
pattern functional, and the field supporting inequality.
Section~\ref{arr:section} supplies the overlap-array and cascade laws.
The remaining stages are as follows.
\begin{enumerate}
 \item \emph{Upper comparison (Section~\ref{upper:section}).}
 Enrich the spins by cascade labels and hierarchical Gaussian fields,
 and Poissonize the number of patterns. If a proposed upper bound fails,
 an affine comparison has a negative minimum in a compact parameter
 region. A quadratic perturbation penalty gives fluctuation control at
 that deterministic minimum, yielding joint Ghirlanda--Guerra identities.
 These order the spin and label overlaps together. Field variations give
 tail inequalities that contradict the density derivative at the minimum.
 \item \emph{Restored covariances (Section~\ref{bulk:section}).}
 Enforce the overlap identities with two patterns omitted, then restore
 them as independent Gaussian marks. The marked-cascade calculation
 evaluates first and second moments of the empirical gradient covariance.
 It becomes a deterministic nondecreasing function of the limiting spin
 overlap; its diagonal and the dilation correction become deterministic
 as well.
 \item \emph{Cavity increments (Section~\ref{cav:section}).}
 Add a fixed number of Ising coordinates. A Taylor expansion and a
 normalized-measure estimate reduce the increment to a Gaussian cavity
 field and finitely many new patterns. Coordinate exchangeability
 identifies the entropy subgradient. Each limiting increment is bounded
 below by the variational value, up to the floor error in the pattern
 count. Averaging the perturbation parameters and telescoping then
 finishes the lower bound.
\end{enumerate}
The array calculations retain independent residual Gaussian noise on every
replica visit, including two visits to the same cascade leaf. This is needed
when the patterns are restored and when cavity-spin magnetizations are
evaluated; Figure~\ref{arr:residual-figure} illustrates the construction.

Finally, Section~\ref{sec:extensions} removes smoothness. Deleting one
pattern gives a Gaussian-$L^1$ approximation bound uniform in dimension;
the tilted recursion gives a matching bound uniform in the overlap path.
These estimates yield the bounded Borel theorem and the separately stated
integrable Gaussian extension. Entrywise replacement proves random-sign
universality under the stated Gaussian-null discontinuity condition.
